\chapter{树、图、BFS 和 DFS}

树和图题的第一难点是建模。题目未必直接给你图，但只要有状态和状态之间的一步转移，就可以看成图。网格题的格子是节点，相邻方向是边；课程表的课程是节点，先修关系是有向边；单词接龙的单词是节点，一次变化是边。

BFS 适合无权图最短步数、层序遍历和多源扩散。腐烂橘子的暴力想法是每一分钟扫描整个网格，看哪些新鲜橘子会腐烂。这样会重复扫描大量空格和已经稳定的格子。优化方法是多源 BFS：初始所有腐烂橘子入队，只从变化边界向外扩散。每个橘子最多入队一次，总复杂度 \complexity{O(mn)}。

\begin{lstlisting}[language=C++,caption={腐烂橘子的 BFS 主循环}]
while (!queue.empty() && fresh_count > 0) {
    const int level_size = static_cast<int>(queue.size());
    for (int step = 0; step < level_size; ++step) {
        const auto [row, col] = queue.front();
        queue.pop();

        for (const auto [dr, dc] : DIRECTIONS) {
            const int next_row = row + dr;
            const int next_col = col + dc;
            if (!in_bounds(next_row, next_col, rows, cols) ||
                grid[next_row][next_col] != FRESH) {
                continue;
            }
            grid[next_row][next_col] = ROTTEN;
            --fresh_count;
            queue.emplace(next_row, next_col);
        }
    }
    ++minutes;
}
\end{lstlisting}

岛屿数量的本质是统计无向图连通块。每个陆地格子是节点，相邻陆地之间有边。每遇到一个未访问陆地，就从它开始 BFS，把整个连通块标记掉，答案加一。

DFS 可以递归写，也可以用显式栈写。递归代码短，但 C++ 在线评测中深度过大可能栈溢出。本书会讲递归思想，但实现上优先强调显式队列、栈和循环。

拓扑排序处理有向依赖。课程表中，入度为 0 的课程表示当前没有前置依赖，可以先学。不断移除入度为 0 的点，如果最后处理了所有课程，说明无环；否则存在环。

本章力扣主线：102 二叉树层序遍历、104 最大深度、98 验证二叉搜索树、236 最近公共祖先、200 岛屿数量、994 腐烂的橘子、207 课程表、210 课程表 II、127 单词接龙。复盘时先写节点是什么、边是什么、访问状态是什么、队列或栈里保存什么。
